\section{第40届复赛第5题：电偶极子场中点电荷运动}
\subsection{题干}
质量为 $m$、电荷量为 $e$ 的正点电荷在固定于原点 $O$ 的电偶极子场中运动，电偶极矩 $\boldsymbol p$ 沿 $z$ 轴正向。设点电荷的位置矢量为 $\boldsymbol r(t)$，其大小为 $r(t)$，与 $\boldsymbol p$ 的夹角为 $\theta(t)$。已知 $r(0)=r_0>0$、$\theta(0)=\theta_0$，真空介电常量为 $\varepsilon_0$。电偶极子的电场和电势分别为
\[
\boldsymbol E=\frac{3(\boldsymbol p\cdot\hat{\boldsymbol r})\hat{\boldsymbol r}-\boldsymbol p}
{4\pi\varepsilon_0r^3},
\qquad
\varphi=\frac{\boldsymbol p\cdot\boldsymbol r}{4\pi\varepsilon_0r^3},
\]
其中 $\hat{\boldsymbol r}$ 为 $\boldsymbol r$ 的单位方向矢量。

（1）假定点电荷在垂直于 $z$ 轴的平面内绕 $z$ 轴作匀速圆周运动，确定 $\theta_0$ 以及圆周运动速度 $v_0$。

（2）假定 $t=0$ 时点电荷静止：（i）求角动量大小 $L(\theta)$；（ii）求径向速率 $v_r(r)$；（iii）求运动有界的条件，并确定边界情形下的竖直平面轨道；（iv）运动无界时求 $r(t)$。
\begin{figure}[htbp]\centering\begin{tikzpicture}[>=Stealth,scale=1.0]
  \fill (0,0) circle (2pt) node[below left] {$O$};
  \draw[->,very thick] (0,0) -- (0,2.05) node[above] {$\boldsymbol p$};
  \draw[->] (0,0) -- (2.2,0) node[right] {$x$};
  \draw[->] (0,0) -- (0,2.2) node[above right] {$z$};
  \draw[thick] (0,0) -- (42:2.15) coordinate (Q);
  \fill (Q) circle (1.7pt) node[above right] {$e,m$};
  \draw[<->,thin] (0.18,0.12) -- ($(Q)+(0.12,0.08)$) node[midway,above] {$r$};
  \draw[->] (0,0.72) arc[start angle=90,end angle=42,radius=.72] node[midway,right] {$\theta$};
  \node[below=0.65cm] at (0.8,-0.2) {原题图5a：电偶极子与点电荷};
\end{tikzpicture}
\caption{第40届复赛第5题原题图5a的TikZ重绘。}\end{figure}
\subsection{参考解答整理}
\paragraph{（1）匀速圆周运动。} $z$ 向合力为零给出 $\cos\theta_0=\pm1/\sqrt3$；向心力方向要求取负根，故
\[\boxed{\cos\theta_0=-\frac1{\sqrt3},\qquad v_0=\sqrt{\frac{3ep}{6\pi\varepsilon_0mr_0^2}}}.\]
\paragraph{（2）（i）角动量随极角的变化。} 角动量定理给出
\[\boxed{L^2=\frac{mep}{2\pi\varepsilon_0}(\cos\theta_0-\cos\theta)}.\]
\paragraph{（2）（ii）径向速率。}
能量守恒进而给出
\[\boxed{v_r^2=\frac{ep\cos\theta_0}{2\pi m\varepsilon_0}\left(\frac1{r_0^2}-\frac1{r^2}\right)}.\]
\paragraph{（2）（iii）有界条件与临界轨道。}
所以 $\theta_0\ge\pi/2$ 时运动有界；临界 $\theta_0=\pi/2$ 时 $r=r_0$，轨道为 $\pi/2\le\theta\le3\pi/2$ 的半圆。
\paragraph{（2）（iv）无界运动的径向距离。}
若 $\theta_0<\pi/2$，则
\[\boxed{r(t)=\sqrt{r_0^2+\frac{ep\cos\theta_0}{2\pi\varepsilon_0mr_0^2}t^2}}.\]
\begin{figure}[htbp]\centering\begin{tikzpicture}[>=Stealth,scale=.9]
  \begin{scope}
    \draw[->] (-2.1,0) -- (2.3,0) node[right] {$x$}; \draw[->] (0,-.2) -- (0,2.2) node[above] {$z$};
    \draw[very thick] (0,0) arc[start angle=0,end angle=180,radius=1.55];
    \fill (0,0) circle(1.4pt) node[below right] {$O$};
    \node[below] at (0,-.75) {$\theta_0=\pi/2$：临界有界半圆};
  \end{scope}
  \begin{scope}[xshift=5.8cm]
    \draw[->] (-.2,0) -- (3.2,0) node[right] {$t$}; \draw[->] (0,-.2) -- (0,2.4) node[above] {$r$};
    \draw[thick,blue!65!black,domain=0:3,samples=80] plot (\x,{sqrt(0.45+0.21*\x*\x)});
    \draw[densely dashed] (0,.67) -- (3,.67) node[right] {$r_0$};
    \node[below] at (1.5,-.5) {$\theta_0<\pi/2$：无界径向运动};
  \end{scope}
\end{tikzpicture}
\caption{依据答案结果绘制的临界半圆与无界径向运动辅助图。}\end{figure}
\subsection{转写依据}
题干依据 \texttt{CPHO40-Q05\_transcription\_v1.md}，答案依据第13至16页及 \texttt{CPHO40-Q05\_formula\_check\_v1.csv}。图5a按原题图重绘；第二图为自绘解析辅助图。本题两问分值为10、30分。
